\documentclass[11pt]{article}
\usepackage[margin=1in]{geometry}
\usepackage{booktabs}
\usepackage{xcolor}
\usepackage{graphicx}
\graphicspath{{../results/}{results/}}
\usepackage{hyperref}
\usepackage{amsmath}
\usepackage{caption}
\usepackage{subcaption}
\usepackage{float}

\definecolor{bestgray}{gray}{0.8}
\newcommand{\best}[1]{\colorbox{bestgray}{\strut\textbf{#1}}}

\title{COMPSCI 589 -- Final Project\\[0.5em]
       \large Benchmarking Classical ML Algorithms on Four Datasets}
\author{Vishal Hanuman \and Akshit}
\date{May 8, 2026}

\begin{document}
\maketitle

% ─────────────────────────────────────────────────────────────────────────────
\section*{Note on Implementations}
% Clearly state whose code was used for each dataset / algorithm.
% e.g., "k-NN and Random Forest applied to Datasets 1--2 use
%         Vishal's implementation from HW1 and HW3 respectively."
% ─────────────────────────────────────────────────────────────────────────────

All machine learning algorithms were implemented from scratch.
No machine learning libraries (e.g., scikit-learn, PyTorch) were used for training or prediction.
\texttt{numpy} and \texttt{matplotlib} were used for numerical computation and plotting.
\texttt{scikit-learn} was used \emph{only} for dataset loading, not for training
or prediction.

\begin{itemize}
  \item \textbf{Datasets 1--2} (Digits, Parkinson's):
        k-NN, Gaussian Naive Bayes, Random Forest, and Neural Network implementations
        by \textbf{Vishal Hanuman} (HW1--HW4).
  \item \textbf{Datasets 3--4} (Rice, Credit Approval):
        Experiments by \textbf{Akshit}. Rice uses k-NN and Decision Tree implementations
        from Akshit's homework code. Credit Approval uses Akshit's Decision Tree and
        Random Forest implementations, with the shared HW4 Neural Network used for the
        additional neural-network comparison.
\end{itemize}

% ─────────────────────────────────────────────────────────────────────────────
\section{Introduction}
% ─────────────────────────────────────────────────────────────────────────────

The goal of this project is to compare the performance of classical machine learning algorithms
on four benchmark datasets by evaluating accuracy and F1-score under stratified 10-fold
cross-validation, with systematic hyperparameter tuning.

% ─────────────────────────────────────────────────────────────────────────────
\section{Datasets}
% ─────────────────────────────────────────────────────────────────────────────

\begin{enumerate}
  \item \textbf{Digits} -- 1,797 instances, 64 numerical features, 10 classes (0--9).
        Loaded using the Digits loader from \texttt{scikit-learn}.
  \item \textbf{Parkinson's Disease} -- 195 instances, 22 numerical features, binary
        (0 = healthy, 1 = Parkinson's).
  \item \textbf{Rice Grains} -- 3,810 instances, 7 numerical features, binary
        (Cammeo vs. Osmancik).
  \item \textbf{Credit Approval} -- 653 instances, 6 numerical + 9 categorical features,
        binary (approved vs. not).
\end{enumerate}

% ─────────────────────────────────────────────────────────────────────────────
\section{Experiments and Analyses}
% ─────────────────────────────────────────────────────────────────────────────

\subsection{Dataset 1 -- Hand-Written Digits}

For the Digits dataset, we evaluated k-NN, Random Forest, and a Neural Network.
All features are numerical pixel intensities, so no feature encoding was required.
k-NN is a natural baseline because digits with similar pixel layouts should be close in
the normalized feature space. Random Forests were tested because they can learn nonlinear
decision rules over pixel thresholds, while the Neural Network was included because it can
combine many pixel-level signals through hidden units. All results below use stratified
10-fold cross-validation and report accuracy and macro F1-score.

\subsubsection*{k-NN Hyperparameter Sweep}

\begin{table}[H]
\centering
\caption{Digits: k-NN performance for different values of $k$.}
\begin{tabular}{rcc}
\toprule
$k$ & Accuracy & Macro F1 \\
\midrule
1  & 0.9872 & 0.9871 \\
3  & 0.9872 & 0.9872 \\
5  & \best{0.9883} & \best{0.9883} \\
7  & 0.9861 & 0.9861 \\
11 & 0.9794 & 0.9793 \\
15 & 0.9806 & 0.9805 \\
21 & 0.9745 & 0.9743 \\
31 & 0.9677 & 0.9674 \\
\bottomrule
\end{tabular}
\end{table}

\subsubsection*{Random Forest Hyperparameter Sweep}

\begin{table}[H]
\centering
\caption{Digits: Random Forest performance for different values of \textit{ntree}.}
\begin{tabular}{rcc}
\toprule
\textit{ntree} & Accuracy & Macro F1 \\
\midrule
5   & 0.8965 & 0.8957 \\
10  & 0.9377 & 0.9373 \\
20  & 0.9605 & 0.9604 \\
30  & 0.9621 & 0.9620 \\
50  & 0.9694 & 0.9695 \\
75  & 0.9716 & 0.9718 \\
100 & 0.9660 & 0.9660 \\
150 & \best{0.9749} & \best{0.9748} \\
\bottomrule
\end{tabular}
\end{table}

\subsubsection*{Neural Network Hyperparameter Sweep}

\begin{table}[H]
\centering
\caption{Digits: Neural Network performance for selected architectures and regularization values.}
\begin{tabular}{lccc}
\toprule
Hidden layers & $\lambda$ & Accuracy & Macro F1 \\
\midrule
\texttt{(32,)}     & 0.0 & 0.9800 & 0.9798 \\
\texttt{(64,)}     & 0.0 & \best{0.9800} & \best{0.9799} \\
\texttt{(32, 16)}  & 0.0 & 0.9756 & 0.9755 \\
\texttt{(64, 32)}  & 0.0 & 0.9761 & 0.9761 \\
\texttt{(32,)}     & 0.1 & 0.9544 & 0.9542 \\
\texttt{(64, 32)}  & 0.1 & 0.9538 & 0.9536 \\
\bottomrule
\end{tabular}
\end{table}

\subsubsection*{Learning Curves}

\begin{figure}[H]
  \centering
  \begin{subfigure}[b]{0.48\textwidth}
    \includegraphics[width=\textwidth]{digits_knn_k_sweep.png}
    \caption{k-NN: Accuracy and F1 vs $k$}
  \end{subfigure}
  \hfill
  \begin{subfigure}[b]{0.48\textwidth}
    \includegraphics[width=\textwidth]{digits_rf_ntree_sweep.png}
    \caption{Random Forest: Accuracy and F1 vs \textit{ntree}}
  \end{subfigure}
  \caption{Hyperparameter curves for Dataset 1 (Digits).}
\end{figure}

\begin{figure}[H]
  \centering
  \begin{subfigure}[b]{0.48\textwidth}
    \includegraphics[width=\textwidth]{digits_nn_config_sweep.png}
    \caption{Neural Network: Accuracy and F1 by configuration}
  \end{subfigure}
  \hfill
  \begin{subfigure}[b]{0.48\textwidth}
    \includegraphics[width=\textwidth]{digits_nn_learning_curve.png}
    \caption{Neural Network: Cost vs training set size}
  \end{subfigure}
  \caption{Neural Network graphs for Dataset 1 (Digits).}
\end{figure}

k-NN achieved the best Digits performance, with $k=5$ producing 0.9883 accuracy and
0.9883 macro F1. Very small values of $k$ already worked well, which suggests that the
normalized pixel space has strong local class structure. Performance declined as $k$
became larger, likely because broad neighborhoods smooth over class boundaries between
visually similar digits.

Random Forest performance improved substantially as the number of trees increased, but it
remained below k-NN and the Neural Network. This is reasonable for image-like data: an
axis-aligned tree split only sees one pixel feature at a time, so the forest needs many
trees to approximate spatial digit patterns. The Neural Network performed strongly, with
the best setting at one hidden layer of 64 units and no regularization. Adding
regularization with $\lambda=0.1$ reduced performance, which indicates that these small
networks were more likely underfitting than overfitting in this experiment. The learning
curve decreased as more training examples were used and then began to flatten, suggesting
that the selected network was converging but had less advantage over the simpler k-NN
classifier on this dataset.

% ─────────────────────────────────────────────────────────────────────────────
\subsection{Dataset 2 -- Parkinson's Disease}

For the Parkinson's dataset, we evaluated k-NN, Gaussian Naive Bayes, Random Forest, and a
Neural Network. The dataset contains 22 numerical voice-measurement features and a binary
target variable, \texttt{Diagnosis}. k-NN tests whether nearby voice measurements identify
similar diagnoses. Naive Bayes provides a simple probabilistic baseline, Random Forests
capture nonlinear feature thresholds, and the Neural Network tests whether a learned
nonlinear representation improves performance. The F1-score reported here is the F1-score
for the positive Parkinson's class.

\subsubsection*{k-NN Hyperparameter Sweep}

\begin{table}[H]
\centering
\caption{Parkinson's: k-NN performance for different values of $k$.}
\begin{tabular}{rcc}
\toprule
$k$ & Accuracy & F1 \\
\midrule
1  & \best{0.9544} & \best{0.9676} \\
3  & 0.9278 & 0.9530 \\
5  & 0.9278 & 0.9528 \\
7  & 0.9325 & 0.9565 \\
11 & 0.8917 & 0.9319 \\
15 & 0.8420 & 0.9034 \\
21 & 0.8373 & 0.9028 \\
31 & 0.8403 & 0.9047 \\
\bottomrule
\end{tabular}
\end{table}

\subsubsection*{Gaussian Naive Bayes Hyperparameter Sweep}

\begin{table}[H]
\centering
\caption{Parkinson's: Gaussian Naive Bayes performance for different variance floors.}
\begin{tabular}{lcc}
\toprule
$\epsilon$ & Accuracy & F1 \\
\midrule
\texttt{1e-9} & 0.6915 & 0.7423 \\
\texttt{1e-7} & 0.6968 & 0.7479 \\
\texttt{1e-5} & 0.6979 & 0.7504 \\
\texttt{0.001} & 0.6929 & 0.7497 \\
\texttt{0.01} & 0.7129 & 0.7651 \\
\texttt{0.1} & \best{0.7171} & \best{0.7695} \\
\bottomrule
\end{tabular}
\end{table}

\subsubsection*{Random Forest Hyperparameter Sweep}

\begin{table}[H]
\centering
\caption{Parkinson's: Random Forest performance for different values of \textit{ntree}.}
\begin{tabular}{rcc}
\toprule
\textit{ntree} & Accuracy & F1 \\
\midrule
5   & 0.8554 & 0.9055 \\
10  & 0.8817 & 0.9245 \\
20  & \best{0.9014} & \best{0.9376} \\
30  & 0.8878 & 0.9286 \\
50  & 0.8920 & 0.9327 \\
75  & 0.8823 & 0.9250 \\
100 & 0.8928 & 0.9324 \\
150 & 0.8911 & 0.9311 \\
\bottomrule
\end{tabular}
\end{table}

\subsubsection*{Neural Network Hyperparameter Sweep}

\begin{table}[H]
\centering
\caption{Parkinson's: Neural Network performance for selected architectures and regularization values.}
\begin{tabular}{lccc}
\toprule
Hidden layers & $\lambda$ & Accuracy & F1 \\
\midrule
\texttt{(8,)}    & 0.0 & 0.8323 & 0.8920 \\
\texttt{(16,)}   & 0.0 & \best{0.8378} & \best{0.8968} \\
\texttt{(8, 4)}  & 0.0 & 0.7648 & 0.8652 \\
\texttt{(16, 8)} & 0.0 & 0.7659 & 0.8627 \\
\texttt{(8,)}    & 0.1 & 0.8303 & 0.8966 \\
\texttt{(16, 8)} & 0.1 & 0.7542 & 0.8598 \\
\bottomrule
\end{tabular}
\end{table}

\subsubsection*{Learning Curves}

\begin{figure}[H]
  \centering
  \begin{subfigure}[b]{0.48\textwidth}
    \includegraphics[width=\textwidth]{parkinsons_knn_k_sweep.png}
    \caption{k-NN: Accuracy and F1 vs $k$}
  \end{subfigure}
  \hfill
  \begin{subfigure}[b]{0.48\textwidth}
    \includegraphics[width=\textwidth]{parkinsons_nb_epsilon_sweep.png}
    \caption{Naive Bayes: Accuracy and F1 vs $\epsilon$}
  \end{subfigure}
  \caption{Hyperparameter curves for k-NN and Naive Bayes on Dataset 2.}
\end{figure}

\begin{figure}[H]
  \centering
  \begin{subfigure}[b]{0.48\textwidth}
    \includegraphics[width=\textwidth]{parkinsons_rf_ntree_sweep.png}
    \caption{Random Forest: Accuracy and F1 vs \textit{ntree}}
  \end{subfigure}
  \hfill
  \begin{subfigure}[b]{0.48\textwidth}
    \includegraphics[width=\textwidth]{parkinsons_nn_config_sweep.png}
    \caption{Neural Network: Accuracy and F1 by configuration}
  \end{subfigure}
  \caption{Random Forest and Neural Network configuration curves for Dataset 2.}
\end{figure}

\begin{figure}[H]
  \centering
  \includegraphics[width=0.62\textwidth]{parkinsons_nn_learning_curve.png}
  \caption{Neural Network learning curve for Dataset 2.}
\end{figure}

k-NN achieved the best Parkinson's performance, with $k=1$ producing 0.9544 accuracy and
0.9676 F1. Larger values of $k$ reduced performance, which suggests that the most useful
signal is very local in the normalized voice-feature space. Because the dataset is small
and imbalanced toward the positive class, larger neighborhoods likely mix less similar
patients and smooth away useful diagnostic differences.

Gaussian Naive Bayes was the weakest model. The voice measurements are likely correlated
with each other, while Naive Bayes assumes conditional independence given the class. The
epsilon sweep shows that increasing the variance floor helped slightly, but it did not
close the gap to the other methods. Random Forest was the second-best method: it handled
nonlinear feature thresholds better than Naive Bayes, but performance peaked at 20 trees
and then varied slightly, which is consistent with the higher variance expected from
bootstrap training on only 195 instances. The Neural Network trailed k-NN and Random
Forest; the deeper architectures were worse, suggesting that the dataset is too small for
the added capacity. Its learning curve generally decreased as more training instances were
used, but the curve was not perfectly smooth because the train/test split is small and
the training procedure is stochastic.

% ─────────────────────────────────────────────────────────────────────────────
\subsection{Dataset 3 -- Rice Grains}

For the Rice Grains dataset, we evaluated k-NN and a Decision Tree. The dataset has
3,810 examples with seven numerical shape descriptors and a binary label
(\texttt{Cammeo} or \texttt{Osmancik}). Since all features are numerical, k-NN can use
normalized distances directly. The Decision Tree provides a comparison based on learned
feature-threshold rules. We used stratified 10-fold cross-validation and reported
accuracy and F1-score.

\subsubsection*{k-NN Hyperparameter Sweep}

\begin{table}[H]
\centering
\caption{Rice: k-NN performance for different values of $k$.}
\begin{tabular}{rcc}
\toprule
$k$ & Accuracy & F1 \\
\midrule
1  & 0.8853 & 0.8994 \\
3  & 0.9118 & 0.9229 \\
5  & 0.9215 & 0.9315 \\
7  & 0.9236 & 0.9333 \\
9  & 0.9239 & 0.9336 \\
11 & \best{0.9249} & \best{0.9345} \\
\bottomrule
\end{tabular}
\end{table}

\subsubsection*{Decision Tree Hyperparameter Sweep}

\begin{table}[H]
\centering
\caption{Rice: Decision Tree performance for different maximum depths.}
\begin{tabular}{rcc}
\toprule
Max depth & Accuracy & F1 \\
\midrule
2  & \best{0.9247} & \best{0.9337} \\
4  & 0.9231 & 0.9329 \\
6  & 0.9134 & 0.9246 \\
8  & 0.9108 & 0.9221 \\
10 & 0.9052 & 0.9175 \\
\bottomrule
\end{tabular}
\end{table}

\subsubsection*{Learning Curves}

\begin{figure}[H]
  \centering
  \includegraphics[width=0.88\textwidth]{rice_knn_dt_sweep.png}
  \caption{Rice: k-NN and Decision Tree hyperparameter sweeps.}
\end{figure}

Both methods performed strongly on Rice. k-NN with $k=11$ achieved the best overall
result, with 0.9249 accuracy and 0.9345 F1. The steady improvement from $k=1$ through
$k=11$ suggests that local neighborhoods are helpful, but the dataset benefits from some
smoothing instead of relying on a single nearest neighbor. The Decision Tree was almost
as strong at depth 2, with 0.9247 accuracy and 0.9337 F1. Deeper trees performed worse,
which suggests that a small number of shape-feature thresholds captures most of the useful
class separation and that additional depth begins to fit fold-specific noise.

% ─────────────────────────────────────────────────────────────────────────────
\subsection{Dataset 4 -- Credit Approval}

For the Credit Approval dataset, we evaluated Decision Tree, Random Forest, and Neural
Network models. This dataset is different from the first three because it mixes numerical
and categorical attributes. The categorical variables were converted into one-hot features
before model training, while numerical variables were kept as continuous values. We again
used stratified 10-fold cross-validation and compared accuracy and F1-score.

\subsubsection*{Decision Tree Hyperparameter Sweep}

\begin{table}[H]
\centering
\caption{Credit Approval: Decision Tree performance for different maximum depths.}
\begin{tabular}{rcc}
\toprule
Max depth & Accuracy & F1 \\
\midrule
2  & \best{0.8636} & \best{0.8627} \\
4  & 0.8255 & 0.8152 \\
6  & 0.8317 & 0.8110 \\
8  & 0.8318 & 0.8161 \\
10 & 0.8181 & 0.7989 \\
15 & 0.8148 & 0.7962 \\
\bottomrule
\end{tabular}
\end{table}

\subsubsection*{Random Forest Hyperparameter Sweep}

\begin{table}[H]
\centering
\caption{Credit Approval: Random Forest performance for different values of \textit{ntree}.}
\begin{tabular}{rcc}
\toprule
\textit{ntree} & Accuracy & F1 \\
\midrule
1  & 0.8086 & 0.7852 \\
5  & 0.8547 & 0.8393 \\
10 & 0.8698 & 0.8586 \\
20 & 0.8650 & 0.8551 \\
30 & \best{0.8728} & \best{0.8637} \\
\bottomrule
\end{tabular}
\end{table}

\subsubsection*{Neural Network Hyperparameter Sweep}

\begin{table}[H]
\centering
\caption{Credit Approval: Neural Network performance for selected architectures and regularization values.}
\begin{tabular}{lccc}
\toprule
Hidden layers & $\lambda$ & Accuracy & F1 \\
\midrule
\texttt{(8,)}    & 0.0 & 0.8608 & 0.8456 \\
\texttt{(16,)}   & 0.0 & 0.8484 & 0.8313 \\
\texttt{(8, 4)}  & 0.0 & 0.8470 & 0.8271 \\
\texttt{(16, 8)} & 0.0 & 0.8547 & 0.8367 \\
\texttt{(8,)}    & 0.1 & \best{0.8623} & \best{0.8552} \\
\texttt{(16, 8)} & 0.1 & 0.5467 & 0.0000 \\
\bottomrule
\end{tabular}
\end{table}

\subsubsection*{Learning Curves}

\begin{figure}[H]
  \centering
  \includegraphics[width=0.88\textwidth]{credit_dt_rf_sweep.png}
  \caption{Credit Approval: Decision Tree and Random Forest hyperparameter sweeps.}
\end{figure}

\begin{figure}[H]
  \centering
  \begin{subfigure}[b]{0.48\textwidth}
    \includegraphics[width=\textwidth]{credit_nn_config_sweep.png}
    \caption{Neural Network: Accuracy and F1 by configuration}
  \end{subfigure}
  \hfill
  \begin{subfigure}[b]{0.48\textwidth}
    \includegraphics[width=\textwidth]{credit_nn_learning_curve.png}
    \caption{Neural Network: Cost vs training set size}
  \end{subfigure}
  \caption{Credit Approval: Neural Network configuration and learning curves.}
\end{figure}

Random Forest gave the strongest Credit Approval result, with \textit{ntree}=30 producing
0.8728 accuracy and 0.8637 F1. The single Decision Tree was close when kept shallow:
depth 2 reached 0.8636 accuracy and 0.8627 F1, but deeper trees declined. That pattern
suggests that the dataset has a few strong approval rules, while deeper splits begin to
memorize smaller subgroups. The Random Forest improved over a single tree by averaging
many bootstrapped trees, which helped stabilize predictions on this mixed-feature dataset.
The Neural Network was competitive but slightly behind the forest. The best network used
one hidden layer of 8 units with $\lambda=0.1$, reaching 0.8623 accuracy and 0.8552 F1.
The larger regularized architecture collapsed to an F1 of 0.0000, which likely means it
predicted mostly one class under that setting. For this dataset, the tree-based models
were a better fit than the deeper neural-network configurations.

% ─────────────────────────────────────────────────────────────────────────────
\section{Summary of Results}
% ─────────────────────────────────────────────────────────────────────────────

% Highlight best result per dataset per metric in gray (\best{value})

\begin{table}[H]
\centering
\caption{Best performance (Accuracy / F1-Score) per algorithm per dataset.
  Gray cells indicate the highest value for each metric on each dataset.}
\scriptsize
\setlength{\tabcolsep}{3pt}
\begin{tabular}{l cc cc cc cc}
\toprule
 & \multicolumn{2}{c}{\textbf{Digits}} & \multicolumn{2}{c}{\textbf{Parkinson's}}
 & \multicolumn{2}{c}{\textbf{Rice}} & \multicolumn{2}{c}{\textbf{Credit}} \\
\cmidrule(lr){2-3}\cmidrule(lr){4-5}\cmidrule(lr){6-7}\cmidrule(lr){8-9}
\textbf{Algorithm} & Acc & F1 & Acc & F1 & Acc & F1 & Acc & F1 \\
\midrule
k-NN            & \best{0.9883} & \best{0.9883} & \best{0.9544} & \best{0.9676} & \best{0.9249} & \best{0.9345} & -- & -- \\
Decision Tree   & -- & -- & -- & -- & 0.9247 & 0.9337 & 0.8636 & 0.8627 \\
Naive Bayes     & -- & -- & 0.7171 & 0.7695 & -- & -- & -- & -- \\
Random Forest   & 0.9749 & 0.9748 & 0.9014 & 0.9376 & -- & -- & \best{0.8728} & \best{0.8637} \\
Neural Network  & 0.9800 & 0.9799 & 0.8378 & 0.8968 & -- & -- & 0.8623 & 0.8552 \\
\bottomrule
\end{tabular}
\end{table}

% ─────────────────────────────────────────────────────────────────────────────
\section{Discussion}
% ─────────────────────────────────────────────────────────────────────────────

Across the four main datasets, no single model dominated every setting, but the patterns
were consistent with the structure of the data. k-NN was strongest on Digits, Parkinson's,
and Rice. Those datasets all have numerical feature spaces where distance is meaningful
after normalization. Digits and Rice especially show that local neighborhoods carry strong
class information.

Tree-based methods were strongest on Credit Approval. That dataset mixes categorical and
numerical features, and the approval label appears to depend on a small set of rule-like
conditions. A shallow Decision Tree already performed well, and the Random Forest improved
slightly by averaging several bootstrapped trees. In contrast, deeper single trees became
worse on both Rice and Credit, which points to overfitting rather than underfitting.

The Neural Network was competitive when there were enough examples and a useful continuous
feature representation, especially on Digits. It was weaker on Parkinson's and Credit.
For Parkinson's, the dataset is very small, so deeper networks add capacity without much
data to support it. For Credit, the one-hot feature representation is sparse and the
tree-based models fit the rule-like structure more naturally. Naive Bayes was the weakest
model where it was tested, which is reasonable because the Parkinson's voice measurements
are likely correlated and violate the model's conditional-independence assumption.

% ─────────────────────────────────────────────────────────────────────────────
\section{\textcolor{red}{Extra Credit}}
% ─────────────────────────────────────────────────────────────────────────────

\subsection*{\textcolor{red}{[EC\#2] Image Classification on Olivetti Faces}}

\textcolor{red}{To satisfy the EC\#2 criterion of selecting a new \emph{challenging}
dataset, we evaluated our k-NN (HW1) and Neural Network (HW4) implementations on the
\textbf{Olivetti Faces} dataset --- 400 grayscale 64x64 images of 40 distinct subjects
(10 images per subject), loaded with a \texttt{scikit-learn} dataset fetcher.
This is an image classification task with 40 classes, far more than any of the four
main project datasets. The full notebook is at
\texttt{notebooks/05\_olivetti\_ec2.ipynb}.}

\textcolor{red}{Random Forest was omitted because the from-scratch implementation is too
slow with 4096 pixel features. k-NN and Neural Network are also strong face-recognition
baselines.}

\begin{table}[H]
\centering
\caption{Olivetti Faces (EC\#2): best result per algorithm, stratified 10-fold CV.}
\begin{tabular}{lccc}
\toprule
Algorithm & Best Hyperparameter & Accuracy & Macro F1 \\
\midrule
k-NN           & $k=1$ & \best{0.9575} & \best{0.9442} \\
Neural Network & \texttt{(128, 64)}, $\lambda=0.0$ & 0.8350 & 0.8011 \\
\bottomrule
\end{tabular}
\end{table}

\begin{figure}[H]
  \centering
  \begin{subfigure}[b]{0.48\textwidth}
    \includegraphics[width=\textwidth]{olivetti_knn_k_sweep.png}
    \caption{k-NN: Accuracy and F1 vs $k$}
  \end{subfigure}
  \hfill
  \begin{subfigure}[b]{0.48\textwidth}
    \includegraphics[width=\textwidth]{olivetti_nn_config_sweep.png}
    \caption{Neural Network: Accuracy and F1 by configuration}
  \end{subfigure}
  \caption{Olivetti Faces (EC\#2): hyperparameter sweeps.}
\end{figure}

\begin{figure}[H]
  \centering
  \includegraphics[width=0.6\textwidth]{olivetti_nn_learning_curve.png}
  \caption{Olivetti Faces (EC\#2): NN learning curve.}
\end{figure}

\textcolor{red}{The Olivetti results support the extra-credit motivation: this is a
small but difficult 40-class image-classification task. k-NN achieved the best result,
with $k=1$ reaching 0.9575 accuracy and 0.9442 macro F1. Performance decreased as $k$
increased, which suggests that each subject's nearest images are highly informative but
larger neighborhoods begin mixing other identities. The Neural Network also learned a
useful classifier, with its best architecture reaching 0.8350 accuracy and 0.8011 macro
F1, but it was weaker than k-NN. This is expected because the dataset has only 400 images
and 40 classes, so the network has limited data for learning a stable face representation
from 4096 pixel-level inputs.}

\end{document}
